%%%%%
%
\section{D\'etermination du type de particules selon la m\'ethode de Korolev et Sussman}

\itemstyle{Op\'erateur~:} {\sf Habit\_2D}
\index{Habit\_2D}

%%%%%OBJET%%%%
\itemstyle{Objet~:}
Cet op�rateur calcule les fractions de types de particules
(Dendrites, aiguilles, sph�res, irr�guli�res) des donn�es d'une sonde PMS 2D-C.

%%%%%ARGUMENTS%%%%
\bigskip
%%%%%ENTREES%%%%
\begin{tabarg}
\multicolumn{3}{l}{\titretab {Entr\'ees}}\\
\hline
$Time_{reg}$ & Vecteur, {\tt Float\_64} & Temps avec fr\'equence de 1 Hertz ($s$)\\
$Image_{2D}$ & Matrice X1024, {\tt Int\_32} & Image\\
$Time_{asyn}$ & Vecteur, {\tt Float\_64} & Temps de la variable d'entr\'ee Image ($s$)\\
$TAS_{2D}$ & Vecteur, {\tt Int\_16} & True Air Speed ($m/s$)\\
$Code\_Hexa\_sonde$ & Coefficient, {\tt String} & Code hexadecimal de la sonde\\
$Masque\_sonde$ & Coefficient, {\tt String} & Masque hexadecimal des donn\'ees images\\
$Bits\_temps$ & Coefficient, {\tt Int\_16} & Nombre de bits de codage du temps dans les s\'eparateurs\\
$Nbval_{as}$ & Coefficient, {\tt Int\_32} & Nombre d'\'el\'ements dans le vecteur de temps ascynchrone\\
$Temps_{debut}$ & Coefficient, {\tt Int\_32} & Temps du d\'ebut de l'analyse en secondes depuis minuit\\
$Temps_{fin}$ & Coefficient, {\tt Int\_32} & Temps de la fin de l'analyse en secondes depuis minuit\\
$Tronq_{debut}$ & Coefficient, {\tt Int\_16} & Num\'ero de la diode \`a consid\'erer (d\'ebut)\\
$Tronq_{fin}$ & Coefficient, {\tt Int\_16} & Num\'ero de la diode \`a consid\'erer (fin)\\
$Freq$ & Coefficient, {\tt Int\_16} & Nombre d'enregistrements \`a traiter ensemble\\
\end{tabarg}

\medskip
%%%%%SORTIES%%%%
\begin{tabarg}
\multicolumn{3}{l}{\titretab {Sorties}}\\
\hline
$Irregulieres_{2D}$ & Vecteur, {\tt Float\_32} & Fraction de particules de forme irr\'eguli\`eres\\
$Dendrites_{2D}$ & Vecteur, {\tt Float\_32} & Fraction de particules en forme de dendrites\\
$Aiguilles_{2D}$ & Vecteur, {\tt Float\_32} & Fraction de particules en forme d'aiguilles\\
$Spheres_{2D}$ & Vecteur, {\tt Float\_32} & Fraction de particules de forme sph\'erique\\
\end{tabarg}


%%%%%FORMULE%%%%
\itemstyle{Formule}

%%% Habit_2D


%%%%%COMMENTAIRE%%%%
\itemstyle{Commentaire~:}
Les donn�es images de la sonde PMS-2D-C sont analys�es en utilisant plusieurs
enregistrements � la fois. Un nombre de 10 semble correct. Pour que l'analyse
(statistique) soit valide, il faut au moins 200 � 400 particules.
Cette m�thode ne calcule la fraction des dendrites que pour les particules
tronqu�es, pour des raisons �videntes. Cette m�thode fonctionne sur l'hypoth�se
que les distributions de ratios des diff�rents types de particules sont
uniques et la distribution totale est une superposition lin�aire des
diff�rentes distributions. Les ratios utilis�s par cette m�thode ont �t�
d�duits � partir de sets de donn�es homog�nes. Les ratios utilis�s sont
diff�rents selon que nous calculons les fractions des particules tronqu�es ou
non-tronqu�es.

Certaines particules sont rejet�es, soient celles qui ont moins de 20 pixels
(impossible de d�tecter le type de particule � cause de la faible r�solution);
de m�me que les particules tronqu�es dont la plus grande partie de la particule
est hors de la zone de donn�es.

Si des valeurs n�gatives sont calcul�es dans les fractions, celles-ci sont
r�gularis�es (m�thode statistique), c'est-�-dire que ces valeurs sont
redistribu�es aux fractions positives de fa�on pond�r�e.

%%%%%AUTEUR%%%%
\itemstyle{Auteur~:}
Christian Pag\'e, CNRM/GMEI/MNP

%%%%%REFERENCE%%%%
\itemstyle{R\'ef\'erences~:}
Korolev-2000 \cite{Korolev-2000}.
